% diagrams/objetos.tex -- las cosas que se cuentan: cada una dibujada
% dentro de la misma caja, de -1 a 1 en x y en y, para que
% tools/gen_numeros.py las pueda poner en fila, en grupos o en el marco
% de diez sin saber qué es cada una (ver OBJETOS en tools/gen_numeros.py,
% que dice qué macro dibuja cada cosa y cómo se llama en singular y en
% plural). Todas las zonas quedan cerradas y rellenas de blanco, para
% que cada una se pueda colorear.
%
% Muchas son los dibujos de «First Words» (diagrams/kit.tex, del
% repositorio de "Aprendo a leer"), encogidos o movidos hasta caber en la
% caja; las demás son nuevas y siguen el mismo estilo de línea.
\tikzset{objeto/.style={line width=1.1pt, line cap=round, line join=round}}

% Una manzana, con el rabito y una hoja.
\newcommand{\objManzana}{%
  \filldraw[fill=white] (0,0.55) .. controls (-0.35,0.9) and (-1.0,0.78) .. (-0.95,0.02)
    .. controls (-0.9,-0.72) and (-0.35,-1.0) .. (0,-0.8)
    .. controls (0.35,-1.0) and (0.9,-0.72) .. (0.95,0.02)
    .. controls (1.0,0.78) and (0.35,0.9) .. (0,0.55) -- cycle;
  \draw (0,0.55) .. controls (0.02,0.8) .. (0.1,0.98);
  \filldraw[fill=white] (0.08,0.85) .. controls (0.32,1.05) and (0.62,0.98) .. (0.7,0.84)
    .. controls (0.5,0.68) and (0.25,0.7) .. (0.08,0.85) -- cycle;}

% Una pelota de playa, con cuatro gajos.
\newcommand{\objPelota}{%
  \filldraw[fill=white] (0,0) circle (0.9);
  \draw (-0.9,0) .. controls (-0.3,0.38) and (0.3,0.38) .. (0.9,0);
  \draw (0,0.9) .. controls (-0.38,0.3) and (-0.38,-0.3) .. (0,-0.9);}

% Un hueso, tumbado: la barra y los cuatro nudos, un solo contorno.
\newcommand{\objHueso}{%
  \filldraw[fill=white] (-0.492,0.16) -- (0.492,0.16)
    arc[start angle=197.2, end angle=-62.7, radius=0.27]
    arc[start angle=62.7, end angle=-197.2, radius=0.27]
    -- (-0.492,-0.16)
    arc[start angle=17.2, end angle=-242.7, radius=0.27]
    arc[start angle=242.7, end angle=-17.2, radius=0.27] -- cycle;}

% Un sol: el círculo y ocho rayos.
\newcommand{\objSol}{%
  \filldraw[fill=white] (0,0) circle (0.5);
  \foreach \a in {0,45,...,315} {\draw (\a:0.66) -- (\a:0.96);}}

% Una huella de Toby.
\newcommand{\objHuella}{%
  \begin{scope}[shift={(0,-0.28)}, scale=1.65] \huellaFW \end{scope}}

% Toby, sentado.
\newcommand{\objToby}{%
  \begin{scope}[shift={(0.38,-0.88)}, scale=0.4] \tobySentadoFW \end{scope}}

% Un globo, con su hilo.
\newcommand{\objGlobo}{%
  \begin{scope}[shift={(0,-0.85)}, scale=0.72] \globoFW{0.2}{-0.2} \end{scope}}

% Una estrella.
\newcommand{\objEstrella}{%
  \begin{scope}[scale=0.95] \estrellaFW \end{scope}}

% Una hoja de otoño.
\newcommand{\objHoja}{%
  \begin{scope}[shift={(0,-0.8)}, scale=1.2] \hojaFW \end{scope}}

% Un corazón.
\newcommand{\objCorazon}{%
  \begin{scope}[shift={(0,-0.82)}, scale=1.35] \corazonFW \end{scope}}

% Un caramelo.
\newcommand{\objCaramelo}{%
  \begin{scope}[scale=1.25] \carameloFW \end{scope}}

% Un pez.
\newcommand{\objPez}{%
  \begin{scope}[shift={(0.2,0)}, scale=0.72] \pezFW \end{scope}}
